\documentclass[11pt,a4paper]{article}
\usepackage[T1]{fontenc}
\usepackage{lmodern}
\usepackage[margin=25mm,headheight=15pt]{geometry}
\usepackage{amsmath,amssymb,amsthm,mathtools}
\usepackage{booktabs,array,tabularx}
\newcolumntype{Y}{>{\raggedright\arraybackslash}X}
\usepackage{microtype}
\usepackage{fancyhdr}
\usepackage[hidelinks]{hyperref}
\hypersetup{pdftitle={SRG-I: Foundational Reconstruction v0.1.1},pdfauthor={Hilmir Frímann Halldórsson},pdfsubject={Minor revision v0.1.1 for GPT closure review; source-grounded mathematical reconstruction}}
\pagestyle{fancy}
\fancyhf{}
\fancyhead[L]{\small SRG-I \;|\; Foundational reconstruction}
\fancyhead[R]{\small v0.1.1 closure review}
\fancyfoot[L]{\small Revision: 10 October 2026 \;|\; Not for publication}
\fancyfoot[R]{\small \thepage}
\renewcommand{\headrulewidth}{0.3pt}
\renewcommand{\footrulewidth}{0pt}
\setlength{\parindent}{0pt}
\setlength{\parskip}{0.4em}
\setlength{\emergencystretch}{1.5em}
\newtheorem{theorem}{Theorem}[section]
\newtheorem{proposition}[theorem]{Proposition}
\newtheorem{corollary}[theorem]{Corollary}
\theoremstyle{definition}
\newtheorem{definition}[theorem]{Definition}
\newtheorem{example}[theorem]{Example}
\newtheorem{remark}[theorem]{Remark}
\newcommand{\R}{\mathbb R}
\newcommand{\C}{\mathbb C}
\newcommand{\phig}{\phi}
\DeclareMathOperator{\sgn}{sgn}
\DeclareMathOperator{\clip}{clip}
\DeclareMathOperator{\Rea}{Re}
\numberwithin{equation}{section}
\renewcommand{\labelitemi}{$\bullet$}

\begin{document}
\pagestyle{fancy}
\thispagestyle{fancy}
{\small\textbf{SRG-I \quad Mathematical source reconstruction}}
\vspace{0.4em}
\begin{center}
{\LARGE\bfseries Symbolic Resonance Geometry\par}
\vspace{0.45em}
{\Large Historical Reconstruction, Recursive Well-Posedness,\par
and Dual-Sector Symmetries\par}
\vspace{0.7em}
{\large Hilmir Fr\'imann Halld\'orsson\par}
\vspace{0.3em}
{\small Revision date: 10 October 2026}
\end{center}
\vspace{0.2em}
\noindent\textit{Revised manuscript, v0.1.1, returned for closure review. Date above is not a publication date.}

\begin{abstract}
The surviving Symbolic Resonance Geometry sources define several state
schemas, operator families and sector relations. We reconstruct their
mathematical content without selecting a replacement evolution. Literal
substitution of a July 2025 temporal second difference into its accompanying
update reduces the scalar problem to $y=C_\lambda(y+b)$. For the printed
golden-ratio threshold scaling we prove the complete solution set for every
real threshold. At positive threshold, some offsets have no successor,
zero offset has a continuum of successors, and a restricted positive branch
has exactly one. The result addresses algebraic existence and uniqueness
for fixed current data, not full well-posedness of the symbolic state.
We distinguish compression maps by carrier, sign, zero and branch behavior,
and separate temporal differences, spatial graph coupling and alternative
toroidal metrics. The June phase pair defines a quarter-turn and real
quadrature; July time shift, geometric reflection and later sector exchange
require different domains and symmetry conditions. Prescribed clock
modulation, trail embeddings and a later driven memory equation likewise
encode different kinds of memory. Proofs and counterexamples identify
recoverable structure while leaving missing state updates, fusion rules,
operator composition and physical interpretations explicit. Later models
serve as dated comparisons, not retrospective completions.
\end{abstract}

\section{Introduction: reconstruction before completion}\label{sec:introduction}

A recursion must specify more than a recurring pattern or a relationship
between quantities. To obtain a next-state map, the state carrier, required
history, external inputs, and update operations must be defined sufficiently
to determine a successor on a stated domain. An implicit relation may instead
have no successor or several. An observer can describe a completed trajectory
without generating one. These distinctions guide the reconstruction below.

The historical corpus uses shared names for constructions that are not the
same mathematical object. The formation manuscript P03 describes persistent
glyph identities, scalar resonance, positions, memory, and echo sets; it also
prints more than one expression for resonance evolution. The numerical
operator manuscript P07 and the later reconstruction P05 use additional,
distinct operator definitions. The June phase constructions P09 and P10
supply a relation between two phase-labelled fields. Their coexistence does
not by itself establish a single composed dynamics
\cite[physical pp.~2--8]{P03} \cite[physical pp.~3--6]{P07}
\cite[physical pp.~1--3]{P05} \cite[physical p.~2]{P09}
\cite[physical p.~2]{P10}.

The question addressed here is therefore not which missing rule should now
be invented, but what the surviving definitions already determine. The
principal result is an exact scalar branch classification for the implicit
relation obtained from P03. The distinction between supported phase identities
and unspecified dynamical symmetries supplies a second part of the analysis.
The present text follows the accepted source reconstruction and branch
appendix, with proof details included rather than replaced by numerical
examples \cite{F01,Branch}.

\subsection{Source and evidence conventions}

Identifiers P01--P19 are stable identifiers within the retained archive, not
publication identifiers. Citations specify physical PDF pages, counted from
one. The formation paper P03 carries a cover date of 14 July 2025; P07 carries
11 July 2025; P09 and P10 carry 29 June 2025. A cover date, a PDF metadata date,
and a date on which a file was copied are different kinds of evidence. The
source catalog retains that distinction \cite{Catalog}.

We separate three roles throughout. A \emph{historical definition} records what
a manuscript or program states. A \emph{deduction} follows from explicitly
stated mathematical assumptions. A \emph{completion choice} supplies additional
structure that the bounded source chain does not determine. A completion
choice can be mathematically useful without being recovered historical intent.
In particular, a scalar interpretation of a symbolic tuple must be declared,
not inferred from notation alone.

The existing Tower II mechanism studies concern a later, specified executable
model. Their completed trajectories do not determine the meaning of an
incomplete 2025 formula. This paper is not a new simulation study and does not
claim that every historical SRG implementation is described by the relation
examined here \cite{F01}.


\subsection{Chronology and limits of the comparison}

The central source chain comprises the June modules and tensor memo,
the 11 July numerical operators, the 13 July signal-sector discussion,
the 14 July formation paper and the 15 July fusion/reconstruction papers.
P11 and P13 have 29 June PDF creation metadata but no established cover
date; this is not evidence of publication priority.
The November geometric and finite-core manuscripts and the March 2026
real coupled-field program are dated comparisons. Their carriers and
operators are not inserted into the July recursion
\cite{P01,P05,P06,P08,P11,P13,P14,P15,Infinity}.
The working title uses \emph{Symbolic Resonance Geometry}; the bibliography
retains each source's actual title rather than asserting a uniform
historical expansion of SRG.

Elementary identities, differentiation, term tests and the contraction
principle are standard tools. The contribution here is their source-specific
application, particularly the exact successor classification and the
separation of incompatible definitions. Priority beyond the bounded corpus
has not been assessed.

\section{State schemas and distinct historical formulas}\label{sec:state}

\subsection{A glyph schema is not yet a closed state space}

P03 defines a glyph by the schema
\begin{equation}
 Q_i=\bigl(\mathrm{ID}_i,R_i(t),\mathbf X_i(t),H_i(t),\mathcal E_i\bigr),
 \label{eq:glyph}
\end{equation}
Here glyph subscripts and the short label $\mathrm{ID}_i$ are presentational
renamings of the source tuple, not extra state laws. The entries describe
persistent identity, resonance, position, memory,
and an echo set. The same manuscript later uses $\Theta_i$ for a symbolic
state that retains resonance, trail direction, and identity
\cite[physical pp.~2,18]{P03}.

Equation~\eqref{eq:glyph} is a useful schema. It does not declare an algebra
in which a scalar can be added to an ID or a set, a norm of the complete tuple,
or updates for all its components. A scalar compressor can act on a selected
numeric component; applying it to the entire tuple requires additional
structure. Consequently, identifying $\Theta_i$ with $R_i$ is a modeling
interpretation, not an equality established by the printed schema.

\subsection{An explicit expression, a consistency expression, and an implicit one}

Three source formulas must be kept distinct. P03 first writes
\begin{equation}
 R_i(t+\Delta t)=C\!\left[R_i(t)+
 \sum_{j\in\mathcal E_i}f\bigl(R_j(t),d_{ij},H_i(t)\bigr)\right]
 \label{eq:explicit}
\end{equation}
and elsewhere gives
\begin{equation}
 R_i(t)=H_i(t)\sum_{j\in\mathcal E_i}f\bigl(R_j(t),d_{ij}\bigr).
 \label{eq:consistent}
\end{equation}
Its trail-based update is
\begin{equation}
 \Theta_i(t+1)=C\!\left[
 \Theta_i(t)+\nabla_T^2\Theta_i(t)+F_i(t)\right].
 \label{eq:implicit-source}
\end{equation}
These appear respectively on physical pages 2, 5, and 7
\cite[\S\S1.1,2.2,3.1]{P03}.

Equation~\eqref{eq:explicit} is first order once its right-hand data and
compressor are supplied. Equation~\eqref{eq:consistent} is an instantaneous
consistency expression. For example, with an empty echo set it forces zero
resonance, while \eqref{eq:explicit} can preserve a nonzero value in an
identity branch of the compressor. Thus they cannot both be imposed on
arbitrary initial amplitudes without a further interpretation. Selecting the
explicit expression as the intended replacement for
\eqref{eq:implicit-source} would be an editorial decision, not a consequence
of the shared notation \cite[\S\S2.1,3.4]{F01}.

We write $F_i$ for the sum of the source's individual feedback terms;
this is shorthand, not an additional force law. Using its printed summand,
the feedback in \eqref{eq:implicit-source} is
\begin{equation}
 F_i(t)=H_i(t)\sum_{j\in\mathcal E_i(t)}
 \frac{R_j(t)\cos\varphi_{ij}(t)}{1+d_{ij}(t)^2}.
 \label{eq:feedback}
\end{equation}
It becomes a scalar expression when its arguments are specified. A complete
model must still state the echo-set selection/update, the distance convention,
the phase or angular-alignment rule, and the memory input. A dimensional
distance also requires a scale in the denominator; choosing a dimensionless
distance is one possible convention, not a derived physical calibration
\cite[physical p.~7, \S3.2]{P03} \cite[\S2.2]{F01}.

Reciprocal echo membership identifies a two-way graph connection. It does not,
without a specified evolution and a state-return condition, prove a periodic
orbit. Similarly, a predicate that two glyphs are close and have similar
resonance does not specify the amplitude, ID, memory, or echo set of a merged
glyph \cite[physical pp.~3,10,12]{P03} \cite[\S\S2.3--2.4]{F01}.

\subsection{Feedback reciprocity, fusion and initialization}
\label{sec:feedback-domain}

The source feedback has useful conditional structure. If distance is
symmetric, $\varphi_{ji}=-\varphi_{ij}$, and echo membership is reciprocal,
then $a_{ij}=\cos\varphi_{ij}/(1+d_{ij}^2)$ is symmetric on the active edges.
The actual coefficient in \eqref{eq:feedback} is $H_i a_{ij}$. Under the
stated symmetry of $a_{ij}$, equality of the two directional coefficients is
equivalent to $(H_i-H_j)a_{ij}=0$. Thus the memory values must agree on
nonzero-weight edges; a zero-weight edge imposes no such condition.
Even when the coefficients are symmetric, there is no compensating term
proportional to $-R_i$. For a mutually coupled pair with positive resonances, aligned phases
and positive memory, both contributions are positive. Symmetric coupling
therefore does not establish total-resonance conservation.
This is a deduction about the printed feedback, not a conservation test on
an executed trajectory \cite[physical p.~7, \S3.2]{P03}.

The corridor criteria include
\begin{equation}
 A\in\mathcal E_B,\quad B\in\mathcal E_A,\quad
 |R_A-R_B|<\epsilon_R,\quad d_{AB}<\lambda_{\rm VP},\quad
 |\varphi_{AB}|<\theta_{\rm lock}.
 \label{eq:corridor}
\end{equation}
The nearby merge statement has the form
\begin{equation}
 d_{AB}<\lambda_{\rm VP},\quad R_A\approx R_B
 \quad\Longrightarrow\quad Q^{\rm fused}_{AB}=\mathrm{ProtoSRG}
 \label{eq:fusion-predicate}
\end{equation}
\cite[physical p.~10, \S\S4.1--4.3]{P03}.
This records a proposed recognition condition. The right side supplies a
name, but not the entries of the tuple \eqref{eq:glyph}, an ID assignment,
a deletion rule for the parents, or an update order. Identifying the
approximation in \eqref{eq:fusion-predicate} with the $\epsilon_R$ tolerance
is plausible but still a declared convention. The cluster condition on
physical p.~12 similarly specifies proximity and amplitude thresholds; it
does not define an invariant set with a basin of attraction.

The fusion paper instead describes an observer pipeline: resonance peaks
are normalized, events $(X,Y,Z,R_{\rm norm})$ are clustered, and events with
the same ID are linked in pass order. Its clustering radius is related to
$\lambda_{\rm VP}$, but the relative scale of spatial coordinates and
normalized amplitude in that distance remains a choice
\cite[physical pp.~5--7, \S\S3.2--3.5]{P01}.
The narrative that several trails reach a common glyph identity is not
an algorithm for merging different IDs. An observer applied to recorded
events and a rule generating those events have different mathematical roles.

P03 calls a pass a computational index rather than continuous physical
time \cite[physical p.~7, \S3.1]{P03}.
The centered temporal expression below requires both a current and a
previous numeric state. A seed at a single pass does not provide the
second time level. No convention for that initialization, nor for the
direction of a stationary trail used in phase alignment, is recovered.
These are explicit domain questions, independent of whether a later
scalar branch equation happens to have a root.


\section{The implicit scalar relation and its complete branches}\label{sec:implicit}

\subsection{Literal reduction}

P03 defines its trail Laplacian by the centered temporal difference
\begin{equation}
 \nabla_T^2\Theta_i(t)=
 \Theta_i(t+1)-2\Theta_i(t)+\Theta_i(t-1).
 \label{eq:temporal}
\end{equation}
No additional coefficient or time-step denominator is printed in this
definition or before the term in \eqref{eq:implicit-source}
\cite[physical p.~8, \S3.3]{P03}.

For the following theorem, interpret the evolving quantity as a finite real
scalar and freeze the current-pass feedback data. Set
\begin{equation}
 y=\Theta_i(t+1),\qquad
 b=\Theta_i(t-1)-\Theta_i(t)+F_i(t).
 \label{eq:offset}
\end{equation}
Substitution, without changing any source coefficient, gives
\begin{equation}
 \boxed{y=C_\lambda(y+b).}
 \label{eq:root}
\end{equation}
The appearance of the future state on both sides is not by itself an error:
an implicit equation can define a valid update. The question is whether this
particular relation has a root and whether that root is unique.

\subsection{Threshold definition and classification}

Write $\phig=(1+\sqrt5)/2$ and $q=\phig^{-1}$. To distinguish the threshold
from other historical parameters, denote it by $\lambda$ throughout this
section. P03's printed threshold compressor is
\begin{equation}
 C_\lambda(x)=
 \begin{cases}
 x,&x<\lambda,\\
 qx,&x\geq\lambda.
 \end{cases}
 \label{eq:compressor}
\end{equation}
Equality belongs to the compressed branch
\cite[physical p.~8, \S3.4]{P03}.

\begin{theorem}[Complete real scalar solution set]
\label{thm:branches}
Let $b,\lambda\in\R$ and let $C_\lambda$ be defined by
\eqref{eq:compressor}. Define
\begin{align}
 \mathcal L(b,\lambda)&=
 \begin{cases}(-\infty,\lambda),&b=0,\\ \varnothing,&b\ne0,\end{cases}
 \\
 \mathcal U(b,\lambda)&=
 \begin{cases}\{\phig b\},&\phig^2b\geq\lambda,\\
 \varnothing,&\phig^2b<\lambda.\end{cases}
\end{align}
Then the complete solution set of \eqref{eq:root} is
\begin{equation}
 \mathcal S(b,\lambda)=\mathcal L(b,\lambda)\cup\mathcal U(b,\lambda).
 \label{eq:solution-set}
\end{equation}
\end{theorem}

\begin{proof}
The input $y+b$ lies in exactly one of the two source branches. If
$y+b<\lambda$, the equation is $y=y+b$. It requires $b=0$, and then every
$y<\lambda$ satisfies both the equation and its branch condition.

If $y+b\geq\lambda$, then $(1-q)y=qb$. The golden-ratio identity
$\phig^2=\phig+1$ gives $q/(1-q)=\phig$, so the only candidate is
$y=\phig b$. Its input is $(\phig+1)b=\phig^2b$; hence this candidate is
admissible exactly when $\phig^2b\geq\lambda$. Both branches have been
exhausted, and each listed value satisfies its own branch condition. This
establishes necessity and sufficiency.
\end{proof}

\begin{corollary}[Positive threshold]
\label{cor:positive}
For $\lambda>0$, the number and form of real scalar successors are as follows.
\end{corollary}

\begin{center}
\renewcommand{\arraystretch}{1.25}
\begin{tabular}{@{}ll@{}}
\toprule
Offset & Complete real solution set \\
\midrule
$b<0$ & $\varnothing$ \\
$b=0$ & $(-\infty,\lambda)$ \\
$0<b<\lambda/\phig^2$ & $\varnothing$ \\
$b\geq\lambda/\phig^2$ & $\{\phig b\}$ \\
\bottomrule
\end{tabular}
\end{center}

\begin{proof}
For a positive threshold, the upper-branch condition forces
$b\geq\lambda/\phig^2>0$. The lower branch exists only at $b=0$.
Equation~\eqref{eq:solution-set} therefore yields the table.
\end{proof}

At the boundary $b=\lambda/\phig^2$, the successor is
$y=\lambda/\phig$ and its input equals $\lambda$, as required by the
source's upper-branch convention. If nonnegative amplitudes are required,
intersect the solution set with $[0,\infty)$. For $\lambda>0$ this changes the
$b=0$ continuum to $[0,\lambda)$ and leaves the other rows unchanged.
At $\lambda=0,b=0$, the real solution set is $(-\infty,0]$, whose nonnegative
part is $\{0\}$. Theorem~\ref{thm:branches} also covers negative thresholds
without extending the positive-threshold table outside its domain
\cite[\S3.2]{F01}.

\subsection{Examples, typing, and what the theorem does not say}

\begin{example}[Failure of universal existence and uniqueness]
Take $\lambda=1$ and set the previous and current scalar states to zero.
Then $b=F$. For $F=0$, both $y=0.2$ and $y=0.7$ are roots, and the entire
interval $(-\infty,1)$ consists of roots. For $F=0.1$, the uncompressed branch
requires $b=0$, while the compressed candidate has input
$0.1\phig^2\approx0.2618<1$ and is inadmissible. There is no root. For $F=1$,
the unique root is $y=\phig$.
\end{example}

Thus the printed relation is not a total, single-valued scalar update over
all finite inputs. Zero feedback defeats universal uniqueness; a small
positive feedback can defeat existence even under a nonnegative-state
restriction. A choice of initial guess or solver tolerance cannot supply a root when
the exact solution set is empty. Conversely, choosing one member of a
continuum of roots requires a stated selection rule.

For finitely many scalar components with all current-pass feedback fixed,
the solution set is the Cartesian product of the componentwise sets.
Existence requires that every factor be nonempty, and uniqueness requires
that every factor be a singleton. If feedback, phase, or echo membership
depends on the unknown successor, this product classification does not solve
the newly coupled problem. Nor does a componentwise interpretation define an
algebra on the tuple \eqref{eq:glyph} \cite[\S3.2]{F01}.

These conclusions concern algebraic existence and uniqueness. They do not
establish continuous dependence, a complete multistep trajectory, stability,
or an attractor. The bounded source investigation did not identify an
implicit solver or branch selector for this relation. The fusion paper's
listing is expressly pseudocode: its resonance-update stub copies and
returns the input, and the text acknowledges the placeholders
\cite[physical pp.~22--24]{P01}. The formation manuscript P03 also contains
the explicit expression \eqref{eq:explicit} \cite[physical p.~2, \S1.1]{P03};
choosing it as a correction to \eqref{eq:root} remains a separate decision
\cite[\S3.4]{F01}.

\subsection{Why the word compression does not settle solvability}

For $\lambda>0$, the one-sided limits of \eqref{eq:compressor} at the threshold
are $\lambda$ and $q\lambda$, so the map has a downward jump. It is not a
globally Lipschitz contraction. Its range is also unbounded above, because
$C_\lambda(x)=qx\to\infty$. Finally, every
$z\in[q\lambda,\lambda)$ has two distinct preimages, $z$ and $z/q$.
These properties follow from the specific printed map; they are not claims
about every operator bearing the name RPCO \cite[\S3.3]{F01}.

\section{Compression and operator identity}\label{sec:compression}

The earlier partner-dependent expression in P03 is
$C(x)=\min(x,qR_{\mathrm{partner}})$. With a fixed partner value and cap
$m$, its substitution into the same implicit structure gives
\begin{equation}
 y=\min(y+b,m),\qquad
 \mathcal S_{\min}(b)=
 \begin{cases}
 \varnothing,&b<0,\\
 (-\infty,m],&b=0,\\
 \{m\},&b>0.
 \end{cases}
\end{equation}
Indeed $y<m$ requires $b=0$, whereas $y=m$ requires $b\geq0$.
This is a different branch classification. If the partner is itself an
unknown next-state value, the fixed-cap argument no longer settles the
coupled system \cite[physical p.~6, \S2.5]{P03} \cite[\S1]{Branch}.

For $L>0$ and $\epsilon>0$, the numerical-operator manuscript P07 supplies
another map:
\begin{equation}
 C_+(x)=\exp\!\left(
 \clip\bigl(\sgn(x)\log(|x|+\epsilon),-L,L\bigr)\right).
 \label{eq:positive-rpco}
\end{equation}
Here the sign is \emph{inside} the exponential. The later signed map places
it outside:
\begin{equation}
 C_s(x)=\sgn(x)\exp\!\left(
 \clip\bigl(\log(|x|+\epsilon),-L,L\bigr)\right).
 \label{eq:signed-rpco}
\end{equation}
The literal expressions imply $C_+(0)=1$ and $C_s(0)=0$. For negative
inputs, the first has positive output while the second preserves the negative
sign. The accepted source comparison records that P07's formula and its
listing agree, despite adjacent sign-preservation prose; changing the
formula to match the prose would be a correction rather than transcription
\cite[physical p.~3]{P07} \cite[\S4.1]{F01}.

Neither logarithmic map should be substituted for
\eqref{eq:compressor} without declaring a different model. Such substitution
also does not automatically restore unique implicit solvability: for the
later defaults, $y=C_s(y)$ has precisely the three roots
$0,\pm e^L$. Appendix~\ref{app:branches} supplies the exhaustive alternate
branch candidates and qualifications \cite[\S\S2--3]{Branch}.

\subsection{Exact sign and zero behavior}
\label{sec:compression-details}

Use $a_-=e^{-L}$ and $a_+=e^L$ for the lower and upper plateaus, and define
$\clip(u,a_-,a_+)=\min(a_+,\max(a_-,u))$. The positive-output expression
\eqref{eq:positive-rpco} is exactly
\begin{equation}
 C_+(x)=
 \begin{cases}
 \clip(x+\epsilon,a_-,a_+),&x>0,\\
 1,&x=0,\\
 \clip((|x|+\epsilon)^{-1},a_-,a_+),&x<0.
 \end{cases}
 \label{eq:positive-branches}
\end{equation}
This follows by monotonicity of the exponential and the identity
$e^{-\log u}=1/u$. At the printed defaults $L=10$ and
$\epsilon=10^{-6}$, for instance, $C_+(-2)=(2+\epsilon)^{-1}>0$.
The adjacent prose says that sign is retained, while the equation and code
both give \eqref{eq:positive-branches}
\cite[physical p.~3, \S2.1]{P07}. Neither version may be silently substituted
for the other on the basis of that prose.

For the later program's $L=9$ and $\epsilon=10^{-9}$,
\begin{equation}
 C_s(x)=
 \begin{cases}
 \sgn(x)\clip(|x|+\epsilon,a_-,a_+),&x\ne0,\\
 0,&x=0.
 \end{cases}
 \label{eq:signed-branches}
\end{equation}
The definition is odd and preserves nonzero signs, but is discontinuous at
zero because the one-sided limits have magnitude
$\max(a_-,\epsilon)>0$. Its plateaus lose magnitude information
\cite[lines 49--53]{Infinity}.
At $b=0$, the relation $y=C_s(y+b)$ has exactly the three roots
$0,\pm a_+$ for these defaults. Appendix~\ref{app:branches} proves this
and lists all sign and plateau candidates without choosing an operator
for the original implicit recurrence.

\subsection{A name does not fix an operator or its carrier}

Other surviving constructions make the distinction especially clear.
The original spatial RPCO helper samples
\begin{equation}
 \rho_{\rm RPCO}(x,y)=
 \sin\!\bigl(5\operatorname{atan2}(y,x)\bigr)
 e^{-0.05(x^2+y^2)}.
 \label{eq:spatial-rpco}
\end{equation}
It generates a field value from position; it is not an update of an incoming
resonance \cite[lines 5--13]{RPCO}.
The velocity-based program instead applies componentwise damping
\begin{equation}
 D_v(v)=\frac{v}{1+\log(1+|v|)}.
 \label{eq:velocity-damping}
\end{equation}
For $v>0$ its derivative is
\[
 \frac{1+\log(1+v)-v/(1+v)}
      {(1+\log(1+v))^2}>0.
\]
The map is odd, strictly increasing, reduces the magnitude of each nonzero
input, and is unbounded as $|v|\to\infty$. Thus damping need not be a cap
\cite[lines 29--31]{Attractor}.
A separate early vector-field program leaves vectors below magnitude 5
unchanged and multiplies the others by 0.2; that threshold filter is also
unbounded \cite[lines 28--32]{Helicity}.

The November finite-core model uses a linear partial projection
$(1-\lambda_c)I+\lambda_c P_{\rm core}$ on a complex vector space
\cite[physical p.~5, \S3.1]{P15}. It neither thresholds a real scalar nor
evaluates \eqref{eq:spatial-rpco}. The neutral organizing principle is the
carrier, domain and literal map. RPCO, TGMO and REFU remain historical names,
but their reuse supplies no equality of operators.

\begin{center}
\small
\renewcommand{\arraystretch}{1.18}
\begin{tabularx}{\textwidth}{@{}p{0.24\textwidth}YY@{}}
\toprule
Source family & Input and operation & Information needed for use\\
\midrule
P03 partner cap & Scalar, partner-dependent minimum & Partner identity and timing\\
P03 threshold & Scalar, identity/scaled branches & Threshold and admissible data\\
P07 logarithmic & Real input, positive output & Regularizer, clamp and composition\\
Original helper & Position to sampled field & Coordinate domain and sampling\\
Velocity program & Velocity to damped velocity & Force and time-step convention\\
Later signed code & Real amplitude to signed clamp & Stage order and parameters\\
P15 finite core & Complex vector, partial projection & Basis, projector and factor order\\
\bottomrule
\end{tabularx}
\end{center}
This table compares definitions; it is not a lineage asserting that one
operation corrected or completed another.


\section{Temporal curvature, spatial graphs and toroidal geometry}
\label{sec:geometry}

\subsection{The index differentiated is part of the definition}

The centered difference \eqref{eq:temporal} acts along pass index and
annihilates affine temporal sequences. A graph difference acts across
simultaneous positions. They may both encode a form of variation, but their
domains and required data differ. The source notation $\nabla_T^2$ does
not erase that distinction.

The reconstruction patch gives an explicit spatial engine
\cite[physical pp.~1--2, \S\S2--3]{P05}. In its surviving code, positions
form an array $X\in\R^{N\times3}$, with one row per glyph. Let
$P\in\R^{N\times N}$ average the initially selected neighbors on each
nonempty row and use the identity row for an empty neighbor set. Write
$\Delta X=(P-I)X$. The operative stages are
\begin{align}
 Y&=X+0.05\Delta X+0.15\tanh(\Delta X),\label{eq:patch-a}\\
 Z&=0.95Y+0.05\overline Y+\xi,\label{eq:patch-b}\\
 X^+&=Z\,\operatorname{diag}
 (1,1,1-0.001\sin(2\pi\,0.244\,s)).
 \label{eq:patch-c}
\end{align}
The mean is broadcast to each position, the hyperbolic tangent acts
componentwise, and $\xi$ denotes the program's Gaussian perturbation.
In \eqref{eq:patch-c}, right multiplication scales coordinate columns, so
only each glyph's third coordinate changes in that stage.
The initial graph is built once. The variables named GAMMA, LAMBDA\_VP
and a computed breathing amplitude do not enter these operative updates
\cite[lines 31--40,45--78]{Position}.
Equations~\eqref{eq:patch-a}--\eqref{eq:patch-c} are a static transcription
of that program, not an execution here.

For a row-stochastic $P$, $(P-I)\mathbf1=0$; constant positions have zero
graph difference. On the other hand,
\[
 \mathbf1^T(P-I)X=0\quad\hbox{for every }X
 \quad\Longleftrightarrow\quad \mathbf1^TP=\mathbf1^T.
\]
Thus unweighted-sum conservation additionally needs column balance. A
directed nearest-neighbor graph need not have it, and subsequent nonlinear
or forcing stages require their own analysis. This conditional identity
does not turn the graph operator into the printed temporal one.

The patch carries positions but not the complete resonance/echo schema
\eqref{eq:glyph}. It supplies a surviving explicit alternative, not a
source-identified solver for \eqref{eq:root}. Choosing temporal curvature
as an observer, a spatial term, or a second-order state equation remains
undecided by their coexistence.

\subsection{What the Laplacian failure analysis establishes}

The June failure note distinguishes a constant radius from a spatial
radius field \cite[physical p.~1, \S\S1--3]{P11}.
A constant has zero spatial derivatives; promoting it to
$r(x,y)=\sqrt{x^2+y^2}$ changes its type and its mathematical behavior.
Off the origin,
\begin{equation}
 r_{xx}=\frac{y^2}{r^3},\qquad
 r_{yy}=\frac{x^2}{r^3},\qquad
 \Delta r=\frac1r.
 \label{eq:radius-laplacian}
\end{equation}
These identities follow by differentiation. At the origin $r$ is not
classically differentiable, so they do not supply a smooth global repair.
The note's reported array-gradient failure is a historical account; no
failure case or old program was rerun for this paper.

There are also distinct spatial formulas. P08's code differentiates the
array named theta twice along two axes, with coefficients $a^{-2}$ and
$r^{-2}$, while its phi argument is unused. It supplies no spacing or
periodic endpoint convention \cite[physical p.~2, \S2.3]{P08}.
P07 prints
\begin{equation}
 \mathcal L f
 =a^{-2}\nabla_\theta^2f+
 b^{-2}\nabla_\varphi^2f+\nabla_r^2f,
 \label{eq:p07-spatial}
\end{equation}
but its listing applies a function named laplacian to theta, phi and r
without a field argument or a definition of that function
\cite[physical p.~4, \S2.2]{P07}.
Differentiating coordinate arrays and differentiating a field with respect
to coordinates are not the same operation. Calling an array gradient
twice also does not impose toroidal boundary conditions.

\subsection{A conditional embedded-torus diagnostic}
\label{sec:torus}

The following calculation specifies its geometry; it is not a recovered
replacement. For an embedded ring torus with major radius $R>r>0$, let
\begin{equation}
 \mathbf x(u,v)=
 ((R+r\cos u)\cos v,(R+r\cos u)\sin v,r\sin u).
 \label{eq:torus-embedding}
\end{equation}
The coordinate tangents are orthogonal with squared lengths
$r^2$ and $(R+r\cos u)^2$. Hence
$g=\operatorname{diag}(r^2,(R+r\cos u)^2)$.
For a smooth scalar field the standard Laplace--Beltrami formula gives
\begin{align}
 \Delta_g f
 &=\frac1{\sqrt{\det g}}\,
   \partial_a\!\left(\sqrt{\det g}\,g^{ab}\partial_bf\right)\nonumber\\
 &=\frac{f_{uu}}{r^2}
   -\frac{\sin u}{r(R+r\cos u)}f_u
   +\frac{f_{vv}}{(R+r\cos u)^2}.
 \label{eq:torus-lb}
\end{align}
The first-order term follows by differentiating the volume factor
$r(R+r\cos u)$; it cannot be discarded while retaining this induced metric.

For $f(u,v)=\cos u$ at $u=\pi/2$,
$\Delta_g f=1/(rR)$, whereas a flat constant-coefficient two-angle
second-derivative operator gives zero there. A flat product torus can
legitimately have the latter type of operator, but it is a different metric.
The comparison distinguishes possible geometries without choosing between
them.

The original TGMO helper samples
\begin{equation}
 ((2+r\cos\theta)\cos(\pi/4),
 (2+r\cos\theta)\sin(\pi/4),r\sin\theta)
 \label{eq:minor-circle}
\end{equation}
\cite[lines 5--14]{TGMO}.
Its fixed azimuth traces one minor circle, repeated by the parameter range.
It specifies neither a full surface field nor a Laplacian.
Wrapped coordinates, a torus-shaped drawing, a flat product metric and
the induced metric of \eqref{eq:torus-embedding} therefore cannot be
identified from nomenclature alone.


\section{Dual-sector relations and meanings of reversal}\label{sec:sectors}
\subsection{The June complex phase pair}

P09 and P10 define
\begin{equation}
 \Phi_Q=e^{i\theta}T,\qquad
 \Phi_D=e^{i(\theta+\pi/2)}T=i\Phi_Q.
 \label{eq:phase-pair}
\end{equation}
This fixes a phase relationship independently of whether a complete update
for $T$ and $\theta$ has been specified
\cite[physical p.~2, \S2.3]{P09} \cite[physical p.~2, \S2.3]{P10}.

\begin{proposition}[Order and orthogonality of the phase relation]
Let $z$ belong to a complex vector space and define $Jz=iz$. Then
$J^2=-I$ and $J^4=I$. If the space has a Hermitian inner product linear in
its second argument, then
\begin{equation}
 \langle z,Jz\rangle=i\lVert z\rVert^2.
\end{equation}
Thus a nonzero pair $(z,Jz)$ is not orthogonal over the complex Hilbert
space. It is orthogonal in the underlying real space equipped with
$(u,v)_{\R}=\Rea\langle u,v\rangle$.
\end{proposition}

\begin{proof}
Successive multiplication by $i$ gives $i^2=-1$ and $i^4=1$.
Linearity in the second inner-product argument gives the displayed identity,
whose real part is zero.
\end{proof}

The resulting sequence
\begin{equation}
 z\ \longmapsto\ iz\ \longmapsto\ -z\ \longmapsto\ -iz
 \ \longmapsto\ z
 \label{eq:quarter-cycle}
\end{equation}
is an algebraic orbit of $J$. It is not evidence that the source dynamics
moves through these four values in consecutive time steps. Likewise, if one
separately adopts a projective-state convention that quotients common phase,
$z$ and $iz$ describe the same ray. The source relation does not supply a
measurement theory identifying distinct particles \cite[\S6.1]{F01}.

Plain exchange of the constrained pair is different. Writing
$E(A,B)=(B,A)$ gives $E^2=I$, but
$E(z,iz)=(iz,z)$ is not of the form $(w,iw)$ unless $z=0$.
Advancing the common phase by $\pi/2$ instead produces $(iz,-z)$, a signed
exchange. Complex conjugation, spatial reflection, and reversal of a
specified evolution are further distinct operations. Their use as a common
word such as ``mirror'' cannot replace a declaration of their actions and
domains \cite[\S6.1]{F01} \cite[\S6]{Branch}.

The July relation $D(t)=Q(t+\tau)$, with $\tau=\pi/(2\omega)$, gives
multiplication by $i$ for the specified mode $Q(t)=e^{i\omega t}$. It does
not do so for an arbitrary signal. A constant is unchanged, while a
$2\omega$ mode acquires a phase of $\pi$. A time shift therefore requires
its own function-space and sampling conventions; it is not automatically a
time-reversal operation \cite[physical pp.~3,8]{P06} \cite[\S6.3]{F01}.

\subsection{Time shift, reflection and reversal}
\label{sec:shift-reversal}

For functions on a two-sided time domain, define
$(S_\tau f)(t)=f(t+\tau)$ and $(Tf)(t)=f(-t)$.
Then
\begin{equation}
 S_\tau^{-1}=S_{-\tau},\qquad T^2=I,\qquad
 TS_\tau T=S_{-\tau}.
 \label{eq:shift-reversal}
\end{equation}
For example, the last identity follows from
$f(-(-t+\tau))=f(t-\tau)$. A shift is invertible but
$S_\tau^2=S_{2\tau}$ is not generally the identity. The July source's
$D(t)=Q(t+\pi/(2\omega))$ therefore relates complete signals; it does not
by itself specify an involution or a causal update on integer passes
\cite[physical pp.~3,8, Eq.~(1) and equation-index entry (1)]{P06}.
For a real monochromatic sinusoid and nonzero $\omega$, quadrature gives
zero inner product when integrated over a full period. The constant-signal
counterexample shows why that orthogonality is not a general property of
the displayed time shift. Its use in a sampled recursion also requires
interpolation, continuation and clock conventions.

A spatial mirror across a plane through $\mathbf x_0$, with specified
unit normal $\mathbf n$, is instead
\begin{equation}
 \mathcal P\mathbf x=
 \mathbf x-2\mathbf n\bigl[\mathbf n\cdot(\mathbf x-\mathbf x_0)\bigr].
 \label{eq:reflection}
\end{equation}
Decomposing $\mathbf x-\mathbf x_0$ into tangential and normal components
proves $\mathcal P^2=I$. The later geometric manuscript also writes
$Q^*=(g(C),H)$, but does not specify the component functions of $g$
\cite[physical p.~10, \S4.2]{P14}.
Neither the star notation nor a reflected tetrahedron establishes
complex conjugation or an action on the original full glyph state.

For an autonomous bijective map $F$, a precise reversal condition is an
involution $\mathcal T$ satisfying
\begin{equation}
 \mathcal T F\mathcal T=F^{-1}.
 \label{eq:reversor}
\end{equation}
This is a definition of reversibility, not a property asserted of the
incomplete source recursion. A nonautonomous map additionally needs the
clock and prescribed forcing reversed; a stochastic model needs a
specified transformation of its noise. Oddness, reciprocity, invertibility
or a spatial reflection alone does not prove \eqref{eq:reversor}.

\subsection{The later real coupled-field formulation}
\label{sec:rd}

The March 2026 program initializes two independent real arrays $R,D$,
with the second draw scaled by 0.01
\cite[lines 70--91]{Infinity}. Its historical labels do not impose
$D=iR$, $D=-R$, or the July time-shift constraint.
These arrays constitute a different state carrier from the June constrained
complex pair.

Freeze the current geometry and let $W$ be its inverse-distance,
row-normalized neighbor matrix. On nonempty rows the simultaneous amplitude
substage is
\begin{equation}
 \begin{pmatrix}R'\\D'\end{pmatrix}
 =B_p\begin{pmatrix}R\\D\end{pmatrix},\qquad
 B_p=
 \begin{pmatrix}
  A_R&K_p\\K_p&A_D
 \end{pmatrix},
 \label{eq:rd-block}
\end{equation}
where
\begin{equation}
 A_R=0.85I+0.15W,\quad A_D=0.925I+0.075W,\quad
 c_p=0.0008699\sin(0.244p+\pi/2).
 \label{eq:rd-parameters}
\end{equation}
Here $K_p=c_pI$ when every row has neighbors. To include empty rows,
set their rows of $W$ to identity rows and the corresponding diagonal
entries of $K_p$ to zero: the program keeps both amplitudes and skips
cross terms there \cite[lines 114--135]{Infinity}.

\begin{proposition}[Exchange and quarter-turn conditions]
\label{prop:rd}
For the fixed-geometry linear substage, sector exchange
$S(R,D)=(D,R)$ commutes with $B_p$ exactly when $A_R=A_D$.
The real quarter-turn $J(R,D)=(-D,R)$ commutes exactly when
$A_R=A_D$ and $K_p=0$.
\end{proposition}
\begin{proof}
Block multiplication gives
\[
 SB_pS=\begin{pmatrix}A_D&K_p\\K_p&A_R\end{pmatrix}.
\]
Equality to $B_p$ is equivalent to equality of the diagonal blocks.
Likewise, comparing
\[
 JB_p=\begin{pmatrix}-K_p&-A_D\\A_R&K_p\end{pmatrix},
 \qquad
 B_pJ=\begin{pmatrix}K_p&-A_R\\A_D&-K_p\end{pmatrix}
\]
gives both asserted conditions.
\end{proof}
For the printed coefficients, equality requires $0.075(W-I)=0$ as
an operator on the chosen domain. On a restricted invariant subspace it
may hold, for example on spatially constant arrays when $W\mathbf1=\mathbf1$.
It does not hold for general nonuniform inputs.

Subsequent program stages use $H(p)$ and $H(p+\pi/2)$ for the two
amplitudes, position fusion weighted by $R$ only, and a sine multiplier
on $R$ only, followed by the signed compressor and position noise
\cite[lines 162--202]{Infinity}.
Consequently, conditional exchange in one substage does not establish
exchange symmetry of the whole update. The cross term in
\eqref{eq:rd-parameters} uses $0.244p$, while the later multiplier uses
$2\pi\,0.244p$; these are different numerical frequencies.

For fixed geometry and fixed clock, the amplitude-only map is odd:
its precompression stages are linear and $C_s$ is odd.
Simultaneous sign inversion is therefore an equivariance of that
conditional map. It fails for the full position/amplitude update:
the center is weighted by $|R|$, but the fusion displacement is multiplied
by signed $R$. Holding positions fixed while reversing amplitudes changes
that displacement.

There is a different supported covariance. Negate all positions, retain
$R,D$, and negate the position-noise realization. Distances and graph
weights stay fixed; componentwise tanh, centers and displacements transform
with the position sign. The full formula therefore transforms by spatial
inversion. A symmetric Gaussian noise law gives invariance in distribution;
using the same nonzero noise realization does not give pathwise equality.
General rotations do not follow, because componentwise tanh is not
equivariant under arbitrary rotations.

Nor does the coupling preserve the whole nonnegative cone. If $c_p<0$,
choose a nonempty row with $R=0$, zero neighboring $R$ and $D>0$.
Its new $R$ is negative. This is an algebraic counterexample for arbitrary
inputs, not a claim about the signs of previously accepted trajectories.
The later program thus supports particular conditional symmetries, not
a general identification of the two sectors.

\subsection{The June tensor expression remains separately typed}
\label{sec:tensor}

The collected June model uses
\begin{equation}
 \kappa(T)=\nabla(\nabla\cdot T)-\Delta T.
 \label{eq:curvature-tensor}
\end{equation}
For a smooth three-dimensional vector field with Euclidean derivatives,
this equals $\nabla\times(\nabla\times T)$ by the standard component
identity. For a scalar in several dimensions the divergence is not
defined by those conventions; higher tensors require specified
contractions. The source does not thereby determine a unique tensor type
\cite[physical pp.~3,5,8--9]{P13}.

For spatially constant phase, linearity gives
$\kappa(e^{i\theta}T)=e^{i\theta}\kappa(T)$.
A spatially varying phase introduces product-rule terms.
The master expression on physical p.~5 places generally complex curvature
inside a floor while projecting only a separate coupling term to its real
part. There is no canonical ordered-real floor on complex numbers.
A real projection or componentwise convention would be additional structure,
not an automatic consequence of the quarter-turn relation.

The same source supplies both a prescribed phase schedule
$\theta_n=n\pi/(2N_{\rm iter})$ and a feedback rule
\begin{equation}
 \theta_{n+1}=\theta_n+
 \epsilon_\theta\,\operatorname{mean}|Q_{\rm RDPTF}|.
 \label{eq:phase-feedback}
\end{equation}
We use $\epsilon_\theta$ to distinguish this coefficient from a logarithmic
regularizer. No rule chooses one schedule as an approximation to the other
\cite[physical pp.~8--9]{P13}.
For $\epsilon_\theta\geq0$ the unwrapped phase is nondecreasing; constant
nonzero input magnitude makes it grow linearly. Reducing phase modulo
$2\pi$ bounds a coordinate, not the tensor amplitude.
These observations retain a well-defined phase identity while leaving the
larger recursion's typing and composition open.


\section{Memory and stability statements with explicit domains}
\label{sec:memory}

\subsection{A prescribed schedule and a stored trail}

Persistent IDs and recorded trails supply one meaning of memory
\cite[physical pp.~2,6]{P01} \cite[physical p.~2]{P03}.
The numerical REFU expression supplies another:
\begin{equation}
 H(t)=2e^{-t/2}\cos t+e^{-0.3t}\sin t+
       1.5e^{-0.7t}+e^{-0.4t}.
 \label{eq:H}
\end{equation}
It depends on the clock and fixed coefficients, with no glyph-history
input \cite[physical p.~6, \S3.1]{P07} \cite[lines 5--13]{REFU}.
Calling it memory does not turn it into an accumulated field record.
The associated factor $1+\lambda H/\max(H)$ needs a normalization domain
and stage placement, which the incomplete composition does not settle
\cite[physical pp.~5--6]{P07}.

For $t\geq0$,
\begin{equation}
 |H(t)|\leq
 2e^{-t/2}+e^{-0.3t}+1.5e^{-0.7t}+e^{-0.4t}\longrightarrow0.
 \label{eq:H-bound}
\end{equation}
The triangle inequality proves the bound. The schedule is not globally
nonnegative: at $t_n=3\pi/2+2\pi n$,
\begin{equation}
 e^{0.3t_n}H(t_n)
 =-1+1.5e^{-0.4t_n}+e^{-0.1t_n}\longrightarrow-1.
 \label{eq:H-sign}
\end{equation}
At integer multiples of $2\pi$ all its nonzero terms are positive.
Thus the clock function has positive and negative values arbitrarily late,
even though both are small. These continuous-clock statements do not
assert signs on every particular discrete sampling schedule.

The later program's shifted factor $H(t+\pi/2)$ changes the decay
coefficients as well as the trigonometric phases: a term $e^{-at}$
acquires the factor $e^{-a\pi/2}$. Different decay rates receive different
factors and the nonoscillating terms remain. It is not a pure quadrature
pair merely because the clock shift is $\pi/2$
\cite[lines 176--182]{Infinity}.

\subsection{Weighted embeddings are observers}

P01's weights and embedding can be written, with $S_j$ denoting the
summed normalized amplitude of trail $j$, as
\begin{equation}
 S_j=\sum_{t\in j}R_{{\rm norm},t},\qquad
 w_j=\frac{S_j}{\sum_kS_k},\qquad E=\sum_jw_jv_j,
 \label{eq:embedding}
\end{equation}
where $v_j$ represents trail information
\cite[physical p.~7, \S3.5]{P01}.
The denominator must be nonzero. If every $S_j\geq0$, these are convex
weights. Likewise, nonnegative $R$ with a positive maximum ensures
$R_{\rm norm}=R/\max R\in[0,1]$; unrestricted signed inputs do not
guarantee this range. These are sufficient sign conditions, not necessary
ones: uniformly negative sums can also normalize to positive weights.

Summation is not injective. With equal weights, the pairs $(-1,1)$ and
$(0,0)$ give the same scalar embedding, and reordering samples does not
change a sum. An embedding can be useful while discarding information;
lossless retention of all ancestral memory does not follow.
The surviving echo observer likewise measures distances between different
positions along one selected glyph's trail, then min--max normalizes the
distance matrix. It is a readout of one trajectory, not a two-glyph feedback
rule; its normalization also needs nonzero range
\cite[lines 10--43]{Echo}.

\subsection{A later state-driven memory equation}

The November geometric model explicitly introduces history through
\begin{equation}
 \dot H(t)=\eta I(t)-\kappa H(t),\qquad
 I(t)=\int_\Omega u(\mathbf r,t)w(\mathbf r)\,d^3r
 \label{eq:memory-ode}
\end{equation}
\cite[physical pp.~5,17, Eqs.~(10),(25)]{P14}.
Suppose $I$ is locally integrable, $\eta,\kappa$ are fixed finite
coefficients and $H(0)$ is specified. Multiplying by $e^{\kappa t}$ and
integrating yields the unique absolutely continuous solution
\begin{equation}
 H(t)=e^{-\kappa t}H(0)
       +\eta\int_0^t e^{-\kappa(t-s)}I(s)\,ds.
 \label{eq:memory-solution}
\end{equation}
Uniqueness follows because the difference of two solutions solves
$\dot h=-\kappa h$, $h(0)=0$.
This is a precise field-history meaning of memory, conditional on the
specified input $u$. It does not solve the coupled field equation or
identify this later $H$ with July's \eqref{eq:H}.

The same source states a wrapped drift--diffusion balance. Integrating its
divergence terms gives zero when the boundary fluxes cancel; total mass is
then conserved if the remaining source has zero spatial integral
\cite[physical p.~17, Eq.~(24)]{P14}.
Boundary cancellation and zero source integral are hypotheses, not
consequences of the word recursion.

\subsection{Positive local results and their limits}

One June scalar module has a complete standard stability argument:
\begin{equation}
 G(x)=\frac{\phig}{\sqrt3}\tanh x,\qquad x\in\R.
 \label{eq:golden-tanh}
\end{equation}
This is the printed $\phig\sqrt3/3$ coefficient
\cite[physical p.~1, \S2.1]{P08}.
Since $|G'(x)|\leq\phig/\sqrt3<1$, the mean-value theorem makes
$G$ a global contraction on the complete real line. Its fixed point
is uniquely zero, which is directly a fixed point. The standard contraction
principle concerns this map, not every composition in the collected model.

A separate drift increment is
$\Delta\Theta_i=\alpha(\Theta_{\rm ref}-\Theta_i)$
\cite[physical p.~13, \S5.4]{P03}.
If this is added to a numeric state with a fixed reference and fixed
$0<\alpha<2$, the new error is $(1-\alpha)$ times the old error, a
strict contraction in any norm. The source does not specify the
reference rule, bounds on its memory-modulated coefficient, or placement
relative to the implicit update. This conditional result is useful
without electing a complete dynamics.

\subsection{A literal series and a geometric contraction claim}

The June collected manuscript prints
\begin{equation}
 \operatorname{Revolution}(X)=
 \sum_{n=1}^{\infty}\Omega_X
 \left(1+\frac{\beta_n}{\alpha_n}\Omega_X\right)
 \label{eq:revolution}
\end{equation}
\cite[physical pp.~4,6]{P13}.
For a fixed finite nonzero $\Omega_X$, every term needs $\alpha_n\ne0$.
Convergence requires its terms to vanish, hence
$\beta_n/\alpha_n\to-1/\Omega_X$. This necessary condition is not
sufficient: writing $\beta_n/\alpha_n=-1/\Omega_X+a_n$ reduces the series
to $\Omega_X^2\sum_n a_n$, which diverges for $a_n=1/n$.
If the ratio tends to zero instead, the terms tend to $\Omega_X\ne0$.
The prose description of a self-modifying constant does not supply an
indexed update for $\Omega_X$ or a convergence prescription. Treating it
as $\Omega_n$ would create an additional definition.

The November field pullback takes the form
\begin{equation}
 u^+=u\circ\Phi,\qquad
 \Phi(\mathbf r)=\mathbf r-\lambda_c\mathbf n(\mathbf r),
 \quad \lambda_c>0.
 \label{eq:normal-map}
\end{equation}
Under its stated assumption that $\Phi$ maps $\Omega$ onto itself,
\[
 \|u\circ\Phi-v\circ\Phi\|_\infty=\|u-v\|_\infty.
\]
This is norm preservation, not strict contraction of fields.
An $L^2$ change of variables instead involves the Jacobian
\cite[physical p.~16, Eqs.~(17)--(18)]{P14}.

There is a separate obstruction to the printed position-space argument.
For a differentiable unit normal,
$\mathbf n^TD\mathbf n=0$, obtained by differentiating
$\mathbf n^T\mathbf n=1$. Therefore $D\mathbf n$ is singular and
$D\Phi=I-\lambda_cD\mathbf n$ has eigenvalue 1. Its induced operator
norm cannot be strictly less than one. A fixed point of the printed
Euclidean-coordinate map would moreover require $\mathbf n=0$,
inconsistent with a unit normal and $\lambda_c>0$.
A vector field vanishing at the target might avoid this obstruction,
but is not the given map.

Even in a different map, eigenvalues inside the unit circle do not alone
imply a contraction in an already specified norm for a nonnormal matrix.
And a contraction on a complete invariant space has one fixed point,
not a nontrivial cycle: if $F^mx=x$ and its contraction factor is $\rho<1$,
then $d(x,Fx)\leq\rho^m d(x,Fx)$, so $Fx=x$.
Thus the source's accompanying alternative of a small set of limit cycles
does not follow from its contraction argument
\cite[physical pp.~16--18]{P14}.
The exact memory solution and conditional mass balance survive these
limitations; a complete attractor theory still needs a defined field map,
domain, boundary conditions and proof.


\section{Later geometric and finite-core comparisons}
\label{sec:later}

Later sources can clarify the range of mathematical meanings attached to
recurring names without completing the earlier model retrospectively.
The November geometric manuscript's elementary identity is
\begin{equation}
 |\sin(3\theta+\delta)+\sin(3\theta-\delta)|
 =2|\sin(3\theta)\cos\delta|.
 \label{eq:sixfold}
\end{equation}
For $\cos\delta\ne0$ the magnitude has period $\pi/3$; for
$\delta=\pi/2$ modulo $\pi$ it vanishes
\cite[physical p.~11, \S4.4]{P14}.
This establishes a sixfold magnitude pattern, not a sixth-harmonic field
or the dynamical emergence of tetrahedra.

Reflection also does not force a Star-of-David projection. Take the
regular tetrahedron with base
\[
 (\cos(2\pi k/3),\sin(2\pi k/3),-1/\sqrt8),\quad k=0,1,2,
\]
and apex $(0,0,3/\sqrt8)$. All six edges have length $\sqrt3$.
Reflecting in $z=0$ reverses its orientation but gives the same XY
projection: a triangle and its center, not two rotated triangles.
This counterexample limits a universal reading of the mirror-projection
claim without ruling out a differently specified geometry
\cite[physical pp.~10--11]{P14}.

The other November manuscript defines a six-component complex model in
the ordered basis $(0+,1+,2+,0-,1-,2-)$
\cite[physical pp.~3--6]{P15}. Following its stated per-corridor
helicity mixing, the order is sector-major and its four factors reduce to
\begin{equation}
 U=A\otimes B,\qquad
 A=e^{-i\alpha\sigma_z}e^{-i\beta\sigma_x},\qquad B=MDC.
 \label{eq:finite-factorization}
\end{equation}
Appendix~\ref{app:finite} specifies the factors and the basis qualification.
A product operator cannot create corridor/helicity entanglement from a
product input. A formula written for a differently ordered basis must
permute all operators consistently before it can be compared.

The same appendix proves the conditional norm identity, parameter
conjugacy and generator-sign comparisons. In particular, explicit
normalization in the source explains normalized reported vectors without
making the underlying linear map unitary
\cite[physical pp.~24,44--45]{P15}.
The source's identification of three spin-half tensor factors with six
dimensions is a separate representation issue: the tensor product has
dimension eight and no total magnetic weight zero. A six-dimensional
direct sum is a different construction
\cite[physical p.~49, Appendix O]{P15}.

These later comparisons are deliberately limited. They do not certify all
later spectral tables, reopen the completed Tower II experiments, or change
the corrected basis and factor-order conventions used by the accepted
QSM/Bridge audits. In particular, a conditional factorization of one
stated per-corridor action is not a claim that every later implementation
uses that action.

\section{Definition requirements and conclusions}
\label{sec:conclusion}

\subsection{What the reconstruction establishes}

The historical materials determine several useful structures. The glyph
schema identifies the intended ingredients of a state. The printed
feedback, compression maps and observer formulas can be compared once
their domains are explicit. The exact branch theorem determines scalar
existence and uniqueness for the implicit temporal relation without
changing its coefficient. The June quarter-turn, July shift, geometric
reflection and later sector matrices support distinct, precisely
conditional identities. The later driven memory equation has an explicit
history representation.

These results support a mathematical reconstruction rather than a selected
replacement dynamics. In particular, Theorem~\ref{thm:branches} concerns
fixed-data scalar solvability, not complete well-posedness of the glyph
system. On its unique positive-threshold branch the root is a linear
function of $b$, but that local formula does not supply roots in the
missing region, choose a member of the $b=0$ continuum, or define updates
for the rest of the tuple.

\subsection{The definitions still needed for a complete dynamics}

A future complete model would have to specify the following mutually
related items. This is a definition checklist, not an implementation plan.

\begin{enumerate}
\item A numeric state carrier, its constraints and norm, and the distinction
between evolving state, persistent labels and readouts.
\item Required time levels and initialization, including whether the
temporal difference generates motion, constrains it, or observes it.
\item A spatial metric, dimensional scales, boundary conventions and any
graph construction or refresh rule.
\item Phase, memory and echo-set laws, including exceptional cases such as
stationary trails, empty neighborhoods and zero normalization denominators.
\item A selected compressor and its stage order, with sign, zero, threshold,
and any partner-timing conventions.
\item A fusion operation if fusion changes state, rather than merely a
predicate applied to observations.
\item For an implicit equation, an admissible domain with existence and
uniqueness conditions, or an explicit multivalued/selection interpretation.
\item A stated symmetry or reversor acting on every relevant state,
parameter, clock and noise component.
\end{enumerate}

None of these can be replaced by the recurrence of an operator name.
Nor can the later source's real $R/D$ arrays supply the missing typing of
the June tensor or July tuple by assumption.

\subsection{Completion comparisons are not recovered intent}

There are several mathematically distinct future choices. One could
formalize the surviving explicit formula \eqref{eq:explicit}; one could
introduce a new implicit model with a globally Lipschitz compressor; or
one could separate temporal curvature from the update and supply a
velocity state. Each needs a declaration and subsequent analysis.
None is selected here.

For example, the standard sufficient condition for a newly specified
equation
\begin{equation}
 y=\widetilde C(\chi y+d),\qquad
 |\chi|\operatorname{Lip}(\widetilde C)<1,
 \label{eq:new-contraction}
\end{equation}
on a complete invariant domain ensures a unique fixed point.
Appendix~\ref{app:new-model} supplies the short proof and the domain
qualification. The original threshold compressor does not satisfy that
global Lipschitz hypothesis. Even introducing $|\chi|<1$ alone leaves
no-root examples because its jump persists. This comparison establishes
what one possible completion would need; it does not approve that completion.

\subsection{Scope of the manuscript}

The analysis is bounded by the retained manuscripts and operative code.
Source equations and counterexamples establish the claims stated here;
they do not authenticate every historical experiment or infer the
author's missing intent. Earlier finite arithmetic checks are reused as
supporting examples, without suite replay. No field update or new
trajectory was computed for the manuscript.

Physical particle labels, horizon narratives and the constant lock are
historical motivations. They need independently defined observables,
units and evidence before they can serve as physical conclusions.
The distinct completed Tower II investigation concerns a specified later
program and remains outside this paper's numerical scope.

The conclusion is therefore constructive but limited: the original
mathematics already supports exact source-specific statements, including
the scalar branch theorem and several conditional symmetry and memory
identities. It does not determine a unique combined evolution. Preserving
that boundary makes those statements available for research without
silently replacing the model whose foundations are being reconstructed.

\appendix
\section{Complete scalar branches for alternative operators}
\label{app:branches}

This appendix compares source-defined operators under substitution into
$y=C(y+b)$. It does not assert that the original paper requested those
substitutions. All variables are finite real numbers unless stated
otherwise. To avoid a collision of roles, the plateau values are denoted
$a_-=e^{-L}$ and $a_+=e^L$, a fixed partner cap is $m$, and a newly
introduced temporal coefficient is $\chi$. These are presentational
renamings, not changes in the source equations.

\subsection{Fixed partner cap}

With the partner held fixed, P03's minimum map gives
$y=\min(y+b,m)$ \cite[physical p.~6, \S2.5]{P03}.
Every root satisfies $y\leq m$.
If $y<m$, equality requires $b=0$; if $y=m$, it requires $b\geq0$.
Consequently the complete set is
\begin{equation}
 \mathcal S_{\min}(b;m)=
 \begin{cases}
 \varnothing,&b<0,\\
 (-\infty,m],&b=0,\\
 \{m\},&b>0.
 \end{cases}
 \label{eq:min-branches}
\end{equation}
Each listed candidate satisfies the original minimum expression, proving
sufficiency. Nonnegative amplitudes require intersection with
$[0,\infty)$; a negative cap then admits no root.
If the partner is an unknown successor, its coupling must be analyzed
separately and this fixed-cap classification is insufficient.

\subsection{Positive-output logarithmic map}
\label{app:positive-log}

Let $L>0$, $\epsilon>0$, $x=y+b$, and use
\eqref{eq:positive-branches}. Every root is positive.
The following exhaustive table includes clamp boundaries in adjacent
branches intentionally; duplicate candidates are removed from the union.

\begin{center}
\small
\renewcommand{\arraystretch}{1.28}
\begin{tabularx}{\textwidth}{@{}p{0.20\textwidth}p{0.27\textwidth}Y@{}}
\toprule
Input branch & Candidate or family & Necessary and sufficient tests\\
\midrule
$x=0$ & $y=1$ & $b=-1$\\
$x>0$, lower & $y=a_-$ & $a_-+b>0$ and $a_-+b+\epsilon\leq a_-$\\
$x>0$, interior & Eligible $y$, only if $b=-\epsilon$ &
 $y-\epsilon>0$ and $a_-\leq y\leq a_+$\\
$x>0$, upper & $y=a_+$ & $a_++b>0$ and $a_++b+\epsilon\geq a_+$\\
$x<0$, lower & $y=a_-$ & $a_-+b<0$ and $(\epsilon-a_--b)^{-1}\leq a_-$\\
$x<0$, interior & $y_\pm$ in \eqref{eq:log-roots} &
 Real discriminant; $y_\pm>0$; $y_\pm+b<0$; $a_-\leq y_\pm\leq a_+$\\
$x<0$, upper & $y=a_+$ & $a_++b<0$ and $(\epsilon-a_+-b)^{-1}\geq a_+$\\
\bottomrule
\end{tabularx}
\end{center}

On the negative-input interior, equality is
\[
 y=\frac1{\epsilon-y-b}
 \quad\Longleftrightarrow\quad y^2+(b-\epsilon)y+1=0,
\]
so the candidates are
\begin{equation}
 y_\pm=\frac{\epsilon-b\pm\sqrt{(b-\epsilon)^2-4}}2.
 \label{eq:log-roots}
\end{equation}
The table's sign and clamp conditions exclude extraneous polynomial
roots. On the positive-input interior, equality is
$y=y+b+\epsilon$, requiring $b=-\epsilon$.
On each plateau the only possible output is its endpoint.
Zero input has the separate value one. These cases exhaust all real
inputs and prove the table.

At the historical values $L=10$, $\epsilon=10^{-6}$ and offset $b=-3$,
both roots of \eqref{eq:log-roots}, approximately $0.381966$ and
$2.618035$, have negative input and lie between the plateaus.
Both solve the substituted relation.
At $b=-\epsilon$, every eligible positive-input interior value is a root.
Thus this alternate compressor does not generally restore unique
solvability. If $L=0$ instead, the map is constant one and $y=1$ is
the unique root for every offset; that degenerate case is outside the
table's assumption.

\subsection{Signed magnitude clamp}

For \eqref{eq:signed-branches}, the complete candidate union is:
\begin{enumerate}
\item $y=0$, admissible exactly when $b=0$.
\item For each $s\in\{-1,+1\}$, the lower plateau $y=sa_-$,
admissible exactly when $s(y+b)>0$ and $|y+b|+\epsilon\leq a_-$.
\item For each $s$, the upper plateau $y=sa_+$, admissible exactly
when $s(y+b)>0$ and $|y+b|+\epsilon\geq a_+$.
\item On the interior with sign $s$, all $y$ satisfying
\[
 b=-s\epsilon,\quad sy>0,\quad s(y+b)>0,\quad
 a_-\leq|y|\leq a_+.
\]
\end{enumerate}
To prove completeness, separate zero and the two input signs.
On a sign-$s$ interior, the equation is $y=y+b+s\epsilon$;
on a plateau the output is its signed endpoint. The listed conditions
are exactly those branch equations and their domains.

For the later defaults $a_->\epsilon>0$ and $b=0$, neither interior
is possible. The lower plateau would require
$a_-+\epsilon\leq a_-$ and fails. Both upper plateaus succeed, as does
zero. There are exactly three roots, $0,\pm a_+$.

\subsection{A newly introduced temporal coefficient}
\label{app:coefficient}

If, as a new comparison only, a coefficient $\chi$ multiplies the
centered temporal term, known current $x$, previous $z$ and feedback $F$
give
\begin{equation}
 y=C_\lambda(\chi y+d),\qquad
 d=(1-2\chi)x+\chi z+F.
 \label{eq:chi}
\end{equation}
The full branch classification, including singular coefficients, is:
\begin{itemize}
\item If $\chi\ne1$, the lower candidate is $d/(1-\chi)$,
valid exactly when that value is less than $\lambda$.
If $\chi=1$, the lower branch exists exactly for $d=0$ and
then contains all $y<\lambda$.
\item If $\chi\ne1/q$, the upper candidate is $qd/(1-q\chi)$,
valid exactly when $d/(1-q\chi)\geq\lambda$.
If $\chi=1/q$, the upper branch exists exactly for $d=0$ and
then contains all $y\geq q\lambda$.
\end{itemize}
Indeed the two branch equations are $y=\chi y+d$ and
$y=q(\chi y+d)$. Solving them and retaining their input inequalities,
with zero denominators handled separately, proves necessity and sufficiency.

For $\chi=1/2,\lambda=1,d=0.4$, only the lower candidate $y=0.8$
is admissible; the upper candidate has input below the threshold.
For $d=0.6$, the lower candidate $1.2$ is inadmissible and the
upper candidate's input is approximately $0.8683<1$, also inadmissible.
There is no root despite $|\chi|<1$. The threshold jump, rather than
merely the coefficient of the implicit term, is consequential.

\subsection{A standard sufficient theorem for a new model}
\label{app:new-model}

Let $E$ be a Banach space and let $\widetilde C:E\to E$ be globally
Lipschitz with constant $K$. For fixed $d\in E$, define
$T(y)=\widetilde C(\chi y+d)$ and suppose $\rho=|\chi|K<1$.
Then
\[
 \|T(y)-T(z)\|\leq\rho\|y-z\|.
\]
The successive differences of the abstract iteration $y_{m+1}=T(y_m)$
are at most $\rho^m\|y_1-y_0\|$. Their geometric-series bound makes
$(y_m)$ Cauchy; completeness gives a limit $y_*$ and continuity gives
$T(y_*)=y_*$. Two fixed points satisfy
$\|y_*-z_*\|\leq\rho\|y_*-z_*\|$ and hence coincide.
This is the standard contraction proof, included to state all needed
assumptions. The same argument works on a closed invariant subset.
Without invariance a local Lipschitz estimate does not establish existence
in the intended domain.

This is a sufficient comparison theorem for a newly specified compressor
and state carrier. It is neither an adopted SRG law nor an executed
iteration. It cannot be applied to the discontinuous threshold map by
omitting that map's branch conditions.


\section{The later finite-core carrier and operator conventions}
\label{app:finite}

\subsection{Basis, action and factorization}

This appendix is conditional on the per-corridor helicity action stated
in the November manuscript. It is not a re-audit of all its tables or of
any corrected kernel. Write the ordered basis as
\[
 (|0,+\rangle,|1,+\rangle,|2,+\rangle,
  |0,-\rangle,|1,-\rangle,|2,-\rangle).
\]
For matrix construction this order identifies the carrier with
$\C^2_{\rm sector}\otimes\C^3_{\rm corridor}$, with the sector index
outermost. Let
\begin{equation}
 s=\frac1{\sqrt3}(1,1,1)^T,\quad P=ss^\dagger,\quad
 C=(1-\lambda_c)I_3+\lambda_cP,
 \label{eq:finite-C}
\end{equation}
\begin{equation}
 M=e^\gamma P+e^{-\eta}(I_3-P),\qquad
 D=\operatorname{diag}(1,e^{-i\varphi},e^{-2i\varphi}).
 \label{eq:finite-MD}
\end{equation}
The full projector is $I_2\otimes P$. In this ordering the four
per-corridor factors are
\begin{align}
 U_{\rm RPCO}&=I_2\otimes C,&
 U_{\rm flip}&=e^{-i\beta\sigma_x}\otimes I_3,\nonumber\\
 U_{\rm TGMO}&=e^{-i\alpha\sigma_z}\otimes D,&
 U_{\rm REFU}&=I_2\otimes M.
 \label{eq:finite-factors}
\end{align}
These express the literal projector, phase and amplification definitions
\cite[physical pp.~3--6, \S\S2--3]{P15}.
With the printed factor order
$U=U_{\rm REFU}U_{\rm TGMO}U_{\rm flip}U_{\rm RPCO}$,
the elementary identity
$(A\otimes B)(C\otimes D)=AC\otimes BD$
proves \eqref{eq:finite-factorization}.
It follows directly that
$U(h\otimes v)=Ah\otimes Bv$, so product inputs remain products.

The explicit index mapping matters:
\begin{center}
\renewcommand{\arraystretch}{1.2}
\begin{tabular}{@{}lll@{}}
\toprule
Convention & Ordered labels & Same-corridor pairs\\
\midrule
Printed sector-major & $0+,1+,2+,0-,1-,2-$ & $(0,3),(1,4),(2,5)$\\
Corridor-major & $0+,0-,1+,1-,2+,2-$ & $(0,1),(2,3),(4,5)$\\
\bottomrule
\end{tabular}
\end{center}
Appendix C writes $I_3\otimes e^{-i\beta\sigma_x}$ beside the
sector-major diagonals \cite[physical p.~24]{P15}.
That Kronecker product mixes adjacent entries, hence implements the
stated per-corridor action in corridor-major order, not the printed order.
The two complete representations are related by a permutation matrix
$\Pi$, with every operator transformed as $\Pi U\Pi^{-1}$.
Changing the flip alone while retaining the other diagonals is not such
a basis change. The factorization above follows the stated action and
declared order; it does not silently certify the inconsistent literal
combination.

\subsection{Norm, normalization and sector conjugacy}

At $\varphi=2\pi/3$, $s^\dagger Ds=(1+e^{-2\pi i/3}
+e^{-4\pi i/3})/3=0$, while $Cs=s$.
For real $\alpha,\beta$, $A$ is unitary. Therefore
\begin{equation}
 \|U(h\otimes s)\|^2=e^{-2\eta}\|h\|^2.
 \label{eq:finite-norm}
\end{equation}
At the source value $\eta=0.423$, the factor is approximately 0.4291.
The raw map does not preserve norm. The source explicitly renormalizes
after each update in Appendix C and again in its implementation discussion
\cite[physical pp.~24,44--45]{P15}.
The map $\psi\mapsto U\psi/\|U\psi\|$ is a nonlinear normalized map
defined where $U\psi\ne0$, rather than a linear unitary operator.

For exchange $S=\sigma_x\otimes I_3$, use
$\sigma_x\sigma_z\sigma_x=-\sigma_z$ and $\sigma_x^2=I$ to obtain
\begin{equation}
 SU(\alpha,\beta)S=U(-\alpha,\beta).
 \label{eq:finite-exchange}
\end{equation}
This is conjugacy between parameter values, not exchange invariance at
the source's fixed $\alpha=2\pi(0.244)$.
Also $e^{-i\beta\sigma_x}=\cos\beta\,I-i\sin\beta\,\sigma_x$
is a complete swap up to a global phase only when
$\beta=\pi/2$ modulo $\pi$. At $\beta=\pi(0.244)$ it is partial mixing
\cite[physical pp.~5--6]{P15}.

For $0\leq\lambda_c<1$, the eigenvalues of $C$ are 1 and
$1-\lambda_c$, so all factors are invertible; at $\lambda_c=1$,
$C$ is singular. Even in the invertible case, reversing the factor order
and inverting amplification/damping does not establish a physical
reversor of the form \eqref{eq:reversor}.

\subsection{Generator convention and finite-step limitations}

Under the convention $U=e^{-iH}$, choosing the memory generator
continuously from zero strength gives
\begin{equation}
 H_{\rm mem}=i[\gamma P_{\rm core}-\eta(I-P_{\rm core})],
 \qquad
 e^{-iH_{\rm mem}}=
 e^{\gamma P_{\rm core}-\eta(I-P_{\rm core})}.
 \label{eq:generator-memory}
\end{equation}
The opposite sign printed on physical pp.~8 and 31 instead inverts
the stated amplification and damping \cite[Appendix E.5]{P15}.
Other logarithm branches require separate choices; the exponential alone
does not give a unique generator.

If $\lambda_c=\kappa_c\Delta t$, expansion of
$I+\kappa_c\Delta t(P_{\rm core}-I)$ gives
$H_{\rm comp}=i\kappa_c(P_{\rm core}-I)$ to first order under
the same convention. This differs both in sign and in its identity term
from the printed $-i\kappa_c(P_{\rm core}-I/3)$
\cite[physical pp.~7,31]{P15}.
For the complete finite operator, $H_{\rm eff}=i\log U$ with a
specified branch; the pseudocode's $-i\log U$ has the opposite sign
\cite[physical p.~34, Appendix H.6]{P15}.
Noncommuting finite factors do not generally have a generator equal to
the sum of their separate generators. A small-step approximation needs
scaled strengths and control of the commutator terms.

\subsection{The representation claim}

The later spin-network identification asserts a tensor product of three
spin-half edge spaces while counting three edges times two helicities as
six states \cite[physical p.~49, Appendix O]{P15}.
A tensor product has dimension $2^3=8$, whereas three doublets in direct
sum have dimension six. In the tensor product the total magnetic weights
are $\pm3/2$ and $\pm1/2$; none is zero. An invariant singlet would be
annihilated by total $J_z$, so no such singlet exists here.
This elementary obstruction concerns that identification, not the
existence of the separate six-component carrier or all possible
representations that one might subsequently define.

\section{Source dictionary, parameter roles and physical interpretations}
\label{app:sources}

The references below give complete source titles and date evidence.
All P-number page locators refer to physical PDF pages. Code locators
refer to retained original lines; the review ledger disambiguates paths,
records hashes, and connects manuscript claims to the page/line evidence.
An archive ID is not a DOI, and no local path is offered as a public link.

\subsection{Several state carriers}
\begin{center}
\small
\renewcommand{\arraystretch}{1.2}
\begin{tabularx}{\textwidth}{@{}p{0.22\textwidth}YY@{}}
\toprule
Source and date & Carrier actually stated & Boundary of the identification\\
\midrule
June P09--P13 & Phase-labelled tensor/field expressions &
Tensor type and composite floor need definition\\
July P03 & Glyph tuple and symbolic trail state &
Tuple schema is not a scalar algebra\\
July P05/code & Positions on a fixed neighbor graph &
No full resonance/echo-state update\\
July P06 & Two signals related by a time shift &
Not arbitrary independent sectors\\
November P14 & Field, geometry and driven memory &
Domains and nonlinear field terms remain required\\
November P15 & Six complex components &
Basis action and factor order must be fixed\\
March Infinity & Two real amplitude arrays and positions &
Not constrained by the June complex relation\\
\bottomrule
\end{tabularx}
\end{center}

\subsection{Constants require a role and an independent estimator}

The fusion and dark-sector manuscripts print
\begin{equation}
 \pi e\phig=\gamma^{-1}\zeta(3)
 \quad\Longleftrightarrow\quad
 \gamma=\frac{\zeta(3)}{\pi e\phig}
 \approx0.0869947475
 \label{eq:constant-lock}
\end{equation}
\cite[physical p.~3, Eq.~(1)]{P01}
\cite[physical p.~4, Eq.~(2)]{P06}.
The latter's symbol appendix also identifies gamma as the
Euler--Mascheroni constant, approximately 0.577
\cite[physical p.~6]{P06}. These are incompatible values under the same
literal identity, not a numerical uncertainty.
The position program's unused GAMMA=0.577, the later finite-core
reinforcement 0.577 and the Infinity coupling parameter 0.08699
play different roles. A shared letter does not derive a relation among them.

The later URF text describes a run-dependent coupling inferred from
spectra, while its algorithm leaves GammaEstimate unspecified
\cite[physical pp.~8--9,13]{P19}.
Solving \eqref{eq:constant-lock} for gamma is not an independent empirical
test of that identity. An estimator would need a definition independent of
the target equality and a declared observable.
Similarly, the source uses of $\lambda_{\rm VP}$ as an overlap length,
clustering radius or numerical 0.618 require different parameter dictionaries.
The number 0.244 cannot silently be radians per pass, cycles per pass and
hertz without a conversion and physical clock.

\subsection{Horizon equations do not specify a reversal map}

The July manuscript's horizon expression is
\begin{equation}
 \Psi_H(t,r)=A_He^{-\Gamma t}\cos(\omega_r r-\varphi_0)
 \label{eq:horizon}
\end{equation}
\cite[physical p.~4, Eq.~(3)]{P06}.
At fixed $r$ it has decay but no temporal cosine. The coefficient
multiplying radius is labeled in hertz in the symbol discussion; a
spatial scale and dimensional convention are missing. The expression
does not supply a scattering boundary condition or a time-reversal action.

For one mode the printed threshold integral evaluates to
\begin{equation}
 \Lambda_{\rm VP}=\int_0^{R_c}\lambda_n e^{-\alpha_n r}\,dr
 =
 \begin{cases}
 \lambda_n(1-e^{-\alpha_nR_c})/\alpha_n,&\alpha_n\ne0,\\
 \lambda_nR_c,&\alpha_n=0 .
 \end{cases}
 \label{eq:threshold-integral}
\end{equation}
No sum over $n$ is printed, so adding one would be a further definition
\cite[physical p.~4, Eq.~(5)]{P06}.
The stated entropy balance similarly implies only
\begin{equation}
 \frac{d}{dt}(S_{\rm child}+S_{\rm parent})=\epsilon_f.
 \label{eq:entropy}
\end{equation}
Conservation needs $\epsilon_f=0$
\cite[physical p.~5, Eq.~(7)]{P06}.
These integral and balance identities can be retained without inferring
microscopic reversed time, or identifying the historical sector labels
with established particles.

\section{Evidence roles and assistance}
\label{app:evidence}

The central statements are supported by the proofs and source locators
in this paper. FOUNDATIONS-01 supplied twenty finite arithmetic
certificates illustrating branch, sign, matrix and constant comparisons
\cite{F01,Branch}. They were retained as prior evidence, not rerun for
a new receipt, and do not substitute for proofs of general claims.
The drafting supplement supplies a claim/source ledger, a cited-page
and code-line review, mathematical notes, build and page-review records,
and consumed-input integrity hashes. It does not reproduce earlier
trajectory archives or certify every numerical claim in the historical
manuscripts.

The proposed human author is Hilmir Fr\'imann Halld\'orsson.
GPT supplied mathematical review of the preceding reconstruction and the
opening draft, then reviewed v0.1 and requested the minor revisions
E01--E03. Codex carried out source investigation, the recorded
FOUNDATIONS-01 algebra checks, manuscript drafting, and the present
revision with citation verification and typesetting review. The prior
algebra checks were not rerun for this revision.
These are assistance roles rather than coauthorship, external peer review,
or publication approval. This v0.1.1 manuscript addresses E01--E03 and is
returned for GPT closure review; acceptance of the revision is not
claimed before that review.
No affiliation, release license or publication venue is assigned.


\begin{thebibliography}{Infinity}
\small
\bibitem[P01]{P01}
Archived manuscript P01.
\emph{Symbolic Recursion and SRG Fusion: Emergent Lattice Fields from Glyph Memory Dynamics}.
Cover date: 15 July 2025. 24 physical PDF pages.
Stable local archive identifier; public publication status not established.

\bibitem[P03]{P03}
Archived manuscript P03.
\emph{Symbolic Recursion and SRG Formation: Full Method and Emergence}.
Cover date: 14 July 2025. 19 physical PDF pages.
Stable local archive identifier; public publication status not established.

\bibitem[P05]{P05}
Archived manuscript P05.
\emph{Patch Research: Detailed Reconstruction of Missing Components in Symbolic Recursion and SRG Framework}.
Cover date: 15 July 2025. 3 physical PDF pages.
Stable local archive identifier; public publication status not established.

\bibitem[P06]{P06}
Archived manuscript P06.
\emph{Recursive Dark Energy and Vacuum Feedback: A Vesica-Phase Cosmological Framework}.
Cover date: 13 July 2025. 9 physical PDF pages.
Stable local archive identifier; public publication status not established.

\bibitem[P07]{P07}
Archived manuscript P07.
\emph{Recursive Numerical Operators for Symbolic Dynamics}.
Cover date: 11 July 2025. 10 physical PDF pages.
Stable local archive identifier; public publication status not established.

\bibitem[P08]{P08}
Archived manuscript P08.
\emph{Recursive Symbolic Evolution and Feedback Dynamics in Python}.
Cover date: 29 June 2025. 3 physical PDF pages.
Stable local archive identifier; public publication status not established.

\bibitem[P09]{P09}
Archived manuscript P09.
\emph{Recursive Symbolic Evolution and Feedback Dynamics in Python}.
Cover date: 29 June 2025. 7 physical PDF pages.
Stable local archive identifier; public publication status not established.

\bibitem[P10]{P10}
Archived manuscript P10.
\emph{Recursive Symbolic Evolution and Its Application in Quantum and Cosmological Systems}.
Cover date: 29 June 2025. 4 physical PDF pages.
Stable local archive identifier; public publication status not established.

\bibitem[P11]{P11}
Archived manuscript P11.
\emph{Why the Toroidal Laplacian Failed: A Structural Explanation}.
No established cover date; PDF creation metadata: 29 June 2025. 2 physical PDF pages.
Stable local archive identifier; public publication status not established.

\bibitem[P13]{P13}
Archived manuscript P13.
\emph{Unified Recursive Cosmological Model Memo}.
No established cover date; PDF creation metadata: 29 June 2025. 17 physical PDF pages.
Stable local archive identifier; public publication status not established.

\bibitem[P14]{P14}
Archived manuscript P14.
\emph{Recursive Engines in Folded Tri-Octagon Geometry: From Symbolic Qubits to the Dual-Tetrahedral SRG Core}.
Cover date: 24 November 2025. 25 physical PDF pages.
Stable local archive identifier; public publication status not established.

\bibitem[P15]{P15}
Archived manuscript P15.
\emph{Recursive SRG Evolution: A Qutrit-Qubit Operator, Dual-Core Geometry, and Computational Interpretation}.
Cover date: 24 November 2025. 51 physical PDF pages.
Stable local archive identifier; public publication status not established.

\bibitem[P19]{P19}
Archived manuscript P19.
\emph{Unified Recursive Fields: Geometry, Recursion, and Consciousness}.
Cover date: 6 September 2025. 15 physical PDF pages.
Stable local archive identifier; public publication status not established.

\bibitem[Infinity]{Infinity}
Archived source code \texttt{srg\_infinity.py}.
Retained March archive source; internal date 14 March 2026. Principal lines: 49--53,70--91,114--135,162--202.
The review ledger records the exact retained path and source identity.

\bibitem[Position]{Position}
Archived source code \texttt{recursive\_field\_evolve.py}.
Original reconstruction engine; no separate composition date established. Principal lines: 31--40,45--78.
The review ledger records the exact retained path and source identity.

\bibitem[RPCO]{RPCO}
Archived source code \texttt{RPCO.py}.
Original Reading material spatial helper. Principal lines: 5--13.
The review ledger records the exact retained path and source identity.

\bibitem[TGMO]{TGMO}
Archived source code \texttt{TGMO.py}.
Original Reading material toroidal sampler. Principal lines: 5--14.
The review ledger records the exact retained path and source identity.

\bibitem[REFU]{REFU}
Archived source code \texttt{REFU.py}.
Original Reading material clock-function helper. Principal lines: 5--13.
The review ledger records the exact retained path and source identity.

\bibitem[Attractor]{Attractor}
Archived source code \texttt{recursive\_attractor.py}.
Original velocity/position program. Principal lines: 29--56.
The review ledger records the exact retained path and source identity.

\bibitem[Helicity]{Helicity}
Archived source code \texttt{helicity\_REFU\_TGMO\_RPCO\_field\_evolution.py}.
Original prescribed vector-field program. Principal lines: 23--40,57--63.
The review ledger records the exact retained path and source identity.

\bibitem[Echo]{Echo}
Archived source code \texttt{echo\_map\_2.py}.
Original \texttt{Simulations\_Trash/Simulations\_part\_4} trail observer. Principal lines: 10--43.
The review ledger records the exact retained path and source identity.

\bibitem[F01]{F01}
\emph{SRG FOUNDATIONS-01 Original Recursion and Dual-Sector Mathematics}.
Codex source-reconstruction report, 10 October 2026, with subsequent GPT review.
Acceptance is limited to the recorded reconstruction scope.

\bibitem[B]{Branch}
\emph{Exact scalar branches and conditional completion theorem}.
Codex, FOUNDATIONS-01 appendix, 10 October 2026.
Analytical comparisons, not an adopted replacement model.

\bibitem[C]{Catalog}
\emph{Source catalog and coverage}.
FOUNDATIONS-01, 10 October 2026. Titles, date evidence, identities and bounded coverage.

\end{thebibliography}
\end{document}
